\subsection{P-15: a width-three obstruction to one centered interpolation pair}
\paragraph{Statement P-15.}
Work in characteristic zero. Put
\[
 F(x_1,x_2,x_3,x_4)
   =2x_1x_2-(x_1+x_2)(x_3+x_4)+2x_3x_4.
 \tag{P15.1}
\]
This nonzero rational polynomial has minimum ROABP width exactly
three, even when the variable order may be chosen freely. It has
a representation in which every matrix factor, and hence every
nonempty interval matrix product, is invertible.

For the four Lagrange nodes
\[
 (a_1,a_2,a_3,a_4)=(1,-1,\mathrm i,-\mathrm i),
 \qquad \mathrm i^2=-1,
 \tag{P15.2}
\]
the standard centered one-pair generator fails:
\[
 F\bigl(t+z\lambda_1(y),\ldots,t+z\lambda_4(y)\bigr)=0.
 \tag{P15.3}
\]
There are also balanced rational-function evaluation failures.
For either ordered set of abscissas
\[
 (1,-1,\mathrm i,-\mathrm i),\qquad (0,1,2,2/3),
 \tag{P15.4}
\]
one has
\[
 F\left(\frac{\alpha a_1+\beta}{\gamma a_1+\delta},
        \ldots,
        \frac{\alpha a_4+\beta}{\gamma a_4+\delta}\right)=0
 \tag{P15.5}
\]
as a rational identity. Thus balanced $\mathrm{RFE}_1^1$ does not
hit this width-three class. For standard integer abscissas
$a_i=i-1$, $n\geq7$, the same failure is witnessed by
\[
 F_{\rm int}
   =2x_1x_4-(x_1+x_4)(x_3+x_7)+2x_3x_7,
 \tag{P15.6}
\]
which uses the harmonic ordered subset $(0,3,2,6)$.

These are distinct coordinate conventions. Equation (P15.3)
concerns the ordinary center $t\mathbf1$ and the four nodes
(P15.2). Equation (P15.6) concerns balanced rational evaluation
at the full standard integer grid. It does not by itself give
a counterexample to the ordinary centered one-pair Lagrange
generator at that integer grid.

\paragraph{Proof: the exact width and rank profile.}
Let $E_{ab}$ be the three-by-three matrix unit. Define
\[
 \begin{aligned}
 M_j(x_j)&=I+u_jx_jE_{12}+v_jx_jE_{23},\\
 (u_1,u_2,u_3,u_4)&=(1,1,-2,0),\\
 (v_1,v_2,v_3,v_4)&=(0,2,-1,-1).
 \end{aligned}
 \tag{P15.7}
\]
Every $M_j$ is upper triangular with diagonal entries one,
so every interval product is invertible. In the $(1,3)$ entry
of the product, a path must take an $E_{12}$ step in an earlier
factor and an $E_{23}$ step in a later factor. Therefore
\[
 e_1^{\mathsf T}M_1M_2M_3M_4e_3
    =\sum_{i<j}u_iv_jx_ix_j=F.
 \tag{P15.8}
\]
This proves width at most three.

For the middle cut $\{1,2\}\mid\{3,4\}$, the left coefficient
rows include the three independent right polynomials
$2x_3x_4$, $-(x_3+x_4)$, and $2$. Their distinct degrees prove
coefficient-matrix rank three, and the fourth possible row is
a repetition. For either other balanced partition, the
degree-one cross block is, up to row and column permutations,
\[
 \begin{pmatrix}2&-1\\-1&2\end{pmatrix},
\]
of determinant three, while the degree-zero and degree-two
blocks both have rank one. Thus those cuts have rank four.
Every order has one of these three balanced partitions at its
middle cut, so every order needs width at least three. The
displayed order has width exactly three.

\paragraph{Proof: ordinary centered interpolation.}
For the nodes (P15.2), Lagrange interpolation gives
\[
 \lambda_a(y)=\frac14
    \left(1+\frac{y}{a}+\frac{y^2}{a^2}+\frac{y^3}{a^3}\right).
 \tag{P15.9}
\]
The polynomial $F$ is invariant under common translation and
homogeneous of degree two. More explicitly, if
$S=x_1+x_2-x_3-x_4$, $U=x_1-x_2$, and $V=x_3-x_4$, then
\[
 2F=S^2-U^2-V^2.
 \tag{P15.10}
\]
At the four interpolants,
\[
 S=y^2,\qquad U=\frac{y^3+y}{2},\qquad
 V=\frac{\mathrm i(y^3-y)}2.
\]
Consequently $S^2=U^2+V^2$, so $F(\lambda(y))=0$.
Translation invariance and homogeneity now give (P15.3).
This is a polynomial identity over $\mathbb Q(\mathrm i)$;
the annihilated polynomial itself has rational coefficients.

\paragraph{Proof: balanced rational evaluation.}
The alternative expression
\[
 F=(x_1-x_3)(x_2-x_4)+(x_1-x_4)(x_2-x_3)
 \tag{P15.11}
\]
makes the rational identity immediate. For
$r(T)=(\alpha T+\beta)/(\gamma T+\delta)$,
\[
 r(a)-r(b)=
 \frac{(\alpha\delta-\beta\gamma)(a-b)}
      {(\gamma a+\delta)(\gamma b+\delta)}.
 \tag{P15.12}
\]
Every ordered quadruple in (P15.4), and also $(0,3,2,6)$,
satisfies the harmonic relation
\[
 (a_1-a_3)(a_2-a_4)+(a_1-a_4)(a_2-a_3)=0.
 \tag{P15.13}
\]
Substitution of (P15.12) into (P15.11) proves (P15.5).
The identity holds on the allowed open set where denominators
are nonzero; singular or constant rational maps cause no
exception. The subset $(0,3,2,6)$ gives (P15.6), with the
other factors in a width-three program taken to be identities.
$\square$

\paragraph{What this closes and what it leaves open.}
The obstruction already occurs with a full-rank matrix word,
so singular rank resets are not needed for failure of one pair.
For these specified node choices, P-14 supplies two pairs for
the full-rank width-three subclass, and (P15.3) proves that one
does not suffice. This disproves a node-uniform one-pair
strengthening of P-14 to width three.

It does not refute balanced $\mathrm{RFE}_{w-1}^{w-1}$ for
general width $w$, nor prove that two pairs suffice for all
width-three programs containing rank-two resets. The broader
width-three construction in I-30 remains valid with six pairs.
The balanced rational ideal and its degree-two generators are
the primary-source framework of Hu, van Melkebeek, and Morgan
(2024), Theorem~1.3, printed p.~6; the identities above are
independently written out and do not assume a general width
theorem from that source. No literature-priority claim is made.
